\documentclass{ctexbook}
\usepackage{Note}
\usepackage{unicode-math}
\setCJKmainfont[
  BoldFont = 思源宋体 Bold,
  ItalicFont = 楷体,
]{宋体}
\setcounter{secnumdepth}{3}
\begin{document}
\chapter{分子对称性}
\section{分子几何构型与对称性}
\subsection{对称操作与对称元素}
为了讨论分子具有的对称性,我们先定义\tbf{对称操作}与\tbf{对称元素}.
\begin{definition}[对称操作与对称元素]
    使一个物体经过几何变换后与变换前完全重合的操作称为\tbf{对称操作}.进行对称操作时依据的几何要素称为\tbf{对称元素}.
\end{definition}
\subsection{常见的对称元素}
\begin{definition}[旋转轴]
    $n$重对称轴$C_n$是一种常见的对称元素,对应的操作是绕轴旋转$\dfrac{2k\pi}{n}$(其中$k=1,2,\cdots,n-1$),分别记作$C_n^1,\,C_n^2,\cdots,\,C_{n}^{n-1}$.
\end{definition}
\begin{definition}[镜面]
    镜面$\sigma$是对称元素,对应的对称操作是沿镜面反射.过主轴的镜面记作$\sigma_{\text{v}}$;垂直于主轴的镜面记作$\sigma_{\text{h}}$;平分两个$C_2$轴的镜面记作$\sigma_{\text{d}}$.
\end{definition}
一般而言,~$\sigma_{\text{d}}$可以视作一种特殊的$\sigma_{\text{v}}$.
\subsection{分子的对称性分类}
为了讨论分子的对称性,将它们按照可能的所有对称操作进行分类.对于单独的分子,所有对称操作至少保持一个点不动,它们构成的集合称作\tbf{点群}.如果在此基础上再考虑平移对称性,则可以进一步分类为\tbf{空间群},这在讨论晶体的对称性时将会介绍.有关群的定义与性质将在下一节介绍.\\
\indent 常见的点群记号有两种,一种是Schoenflies记号,例如$D_{3\text{h}}$, $C_{2\text{v}}$等等;另一种是Hermann-Mauguin记号(又称作国际记号),例如$4\,\text{m}\,\text{m}$, $2/\text{m}$等.描述分子点群时常使用Schoenflies记号.下面是主要的分子点群的分类:
\begin{enumerate}[label=\tbf{\roman*.}]
    \item $C_1$, $C_\i$和$C_\text{s}$群.\\
    $C_1$群只含有恒等操作$E$; $C_\i$群在$C_1$群的基础上含有反演中心$\i$; $C_{\text{s}}$群在$C_1$群的基础上含有镜面$\sigma$.
    \item $C_{n}$, $C_{n\text{v}}$和$C_{n\text{h}}$群.\\
    所有$C_n$群在$C_1$群的基础上含有一根$C_n$旋转轴.此外, $C_{n\text{v}}$点群有$n$个$\sigma_{\text{v}}$镜面; $C_{n\text{h}}$点群有$\sigma_{\text{h}}$镜面.
    \item $D_{n}$, $D_{n\text{h}}$和$D_{n\text{d}}$群.\\
    $D_n$群在$C_n$群的基础上含有$n$个垂直于主轴的$C_2$轴,记作$C_2'$. $D_{n\text{h}}$点群在$D_n$点群的基础上有$\sigma_{\text{h}}$镜面.$D_{n\text{d}}$点群在$D_n$点群的基础上有$n$个$\sigma_{\text{d}}$镜面.
    \item $C_{\infty\text{v}}$和$D_{\infty\text{h}}$群.\\
    线形分子按有无对称中心分属$D_{\infty\text{h}}$和$C_{\infty\text{v}}$点群.
    \item $S_n$群.\\
    $S_n$群在$C_1$群的基础上有$n$重反轴$S_n$.
    \item 含有多个主轴的群.
\end{enumerate}
\subsection{对称性的应用}
\begin{theorem}[用对称性判断偶极矩]
    对称操作不能改变偶极矩,因此只有$C_n$, $C_{n\text{v}}$和$C_{\text{s}}$群可能有偶极矩.
\end{theorem}
\begin{theorem}[用对称性判断手性]
    手性分子不具有虚操作元素(对称中心$\i$,镜面$\sigma$和映轴$S_n$).
\end{theorem}
\section{群论}
\subsection{群的定义}
\begin{definition}[群]
    如果集合$G$中的元素和其上定义的乘法操作满足
    \begin{enumerate}[label=\tbf{\alph*.}]
        \item 对乘法封闭,即
        \[\forall a,b\in G,\quad a\cdot b\in G\]
        \item 满足乘法结合律,即
        \[\forall a,b,c\in G,\quad (a\cdot b)\cdot c=a\cdot(b\cdot c)\]
        \item 有唯一单位元$e$,使得
        \[\forall a\in G,\quad a\cdot e=e\cdot a=a\]
        \item 任意元素都存在唯一逆元,即
        \[\forall a\in G,\quad \exists!b\in G,\quad a\cdot b=b\cdot a=e\]
        可以记作$b=a^{-1}$.
    \end{enumerate}
\end{definition}
我们主要研究的是\tbf{点群}.
\begin{definition}[点群]
    通过同一点的对称操作构成的群称作点群.
\end{definition}
\begin{definition}[类]
    对于群$G$中的元素$a$和$b$,如果存在$g\in G$使得$gag^{-1}=b$,则称$a$与$b$共轭.群$G$中所有与$a$共轭的元素称作$a$的共轭类/等价类.
\end{definition}
\begin{definition}[同态与同构]
    对于群$G$和$H$,如果存在映射$h:G\to H$使得对任意$u,v\in G$都有$h(uv)=h(u)h(v)$,则称$G$与$H$\tbf{同态}.如果$h$是双射,则称$G$与$H$\tbf{同构}.
\end{definition}
\subsection{群表示论}
\begin{definition}[群表示]
    群$G$在向量空间$V$上的\tbf{表示}是一个同态映射
    \[\Gamma:G\mapsto\text{GL}(V)\]
    其中$\text{GL}(V)$是$V$上的线性变换群.
\end{definition}
如此,可以将群元$g$(对于点群而言就是各种对称操作)表示为线性变换,进而表示为矩阵$\Gamma(g)$.
\begin{definition}[可约表示与不可约表示]
    如果群$G$的群元$g$在向量空间$V$上的表示$\Gamma(g)$没有非平凡不变子空间,则称$\Gamma$是一个\tbf{不可约表示}.反之,如果存在非平凡不变子空间,则称$\Gamma$是一个\tbf{可约表示}.
\end{definition}
可约表示可以被约化为不可约表示的直和.具体而言,将向量空间$V$做直和分解,使得群元$g$的表示$\Gamma(g)$满足
\[\Gamma(g)=\bigoplus_{i=1}^{m}\Gamma^{(i)}(g)=\begin{bmatrix}
\Gamma^{(1)}(g) & 0 & \cdots & 0\\
0 & \Gamma^{(2)}(g) & \cdots & 0\\
0 & 0 & \ddots & 0\\
0 & 0 & \cdots & \Gamma^{(m)}(g)
\end{bmatrix}\]
其中每个$\Gamma^{(i)}(g)$都是不可约表示,也就是将$\Gamma(g)$表示为分块对角矩阵,每个块对应一个不可约表示,其作用的空间是$V$的子空间.有关不可约表示的数目以及维数有以下两个定理:
\begin{theorem}[不可约表示的数目]
    群$G$的不可约表示的数目等于群$G$的共轭类的数目.
\end{theorem}
\begin{theorem}[不可约表示的维数]
    群$G$的所有不可约表示的维数的平方和等于$G$的阶(即群元的数目),即
    \[\sum_{i=1}^{m}d_i^2=h\]
\end{theorem}
\subsection{特征标与特征标表}
\begin{definition}[特征标]
    考虑群$G$在特定的向量空间$V$上的表示$\rho$,群元$g$在$V$上的表示$\Gamma(g)$的\tbf{特征标}$\chi(g)$定义为
    \[\chi^{\Gamma}(g)=\text{Tr}(\Gamma(g))\]
\end{definition}
群元$g$及其共轭类的特征标都相同,因为矩阵的迹在相似变换下保持不变.\\
\indent 考虑将群$G$到向量空间$V$上的表示$\Gamma$分解为不可约表示
\[\Gamma=\bigoplus_{i=1}^{m}\Gamma^{(m)}\]
对于每个群元$g$,都可以计算$\Gamma^{\Gamma^{(m)}}(g)$的特征标$\chi^{(m)}(g)$,如此可以构成一张表,每行对应一个不可约表示,每列对应一个共轭类,表中的元素就是对应群元的特征标.这张表称作\tbf{特征标表}.
\begin{theorem}[正交归一定理]
    特征标表的任意不同两行正交,每行自身归一,即
    \[\dfrac{1}{h}\sum_{C}N(C)\chi^{\Gamma^{(i)}}(C)\chi^{\Gamma^{(j)}}(C)=\delta_{ij}\]
    其中$C$表示共轭类.
\end{theorem}
\begin{theorem}[不可约表示的标记]
    通常用$\text{A}$和$\text{B}$标记一维不可约表示,前者对主轴旋转保持不变,即$\chi^{\text{A}}(C_n)=1$,后者对主轴旋转改变符号,即$\chi^{\text{B}}(C_n)=-1$.\\
    通常用$\text{E}$标记二维不可约表示,用$\text{T}$标记三维不可约表示,等等.
\end{theorem}
\section{对称性的应用}
\subsection{消失积分}
对于区域$\Omega$上函数$f$的积分,在简单的情形下可以根据奇偶性等简单对称性判断积分是否为零.对于复杂的情形,可以用群论的方式判断.
\begin{theorem}[用群论判断消失积分]
    对于函数$f$在空间$\Omega$上的积分$\displaystyle\int_{\Omega}f\d\tau$,当且仅当被积函数$f$生成$\Omega$所属点群$G$的全对称不可约表示时,积分不为零.
\end{theorem}
全对称不可约表示的特征标均为$1$,典型的是$\text{A}_1$不可约表示.\\
\indent 我们进一步考虑被积函数是两个函数$f_1$与$f_2$的乘积,即
\[I=\int f_1f_2\d\tau\]
假定$f_1$和$f_2$分别生成表示$\Gamma^{(1)}$和$\Gamma^{(2)}$,根据同态映射的性质可知$f_1f_2$生成表示$\Gamma^{(1)}\otimes\Gamma^{(2)}$.因此,当且仅当$\Gamma^{(1)}\otimes\Gamma^{(2)}$包含全对称不可约表示时,积分$I$不为零.表示的直积的特征标满足
\[\chi^{\Gamma^{(1)}\otimes\Gamma^{(2)}}(g)=\chi^{\Gamma^{(1)}}(g)\chi^{\Gamma^{(2)}}(g)\]
\begin{theorem}[广义正交定理]
    不可约表示$\Gamma$在某表示中出现的次数$n(\Gamma)$可由
    \[n(\Gamma)=\dfrac{1}{h}\sum_{C}N(C)\chi^{\Gamma}(C)\chi(C)\]
    得到,其中$\chi^{\Gamma}(C)$和$\chi(C)$分别表示$\Gamma$和待分解表示的特征标.
\end{theorem}
\begin{theorem}[不可约表示的直积与全对称表示的关系]
    不可约表示$\Gamma^{(i)}$和$\Gamma^{(j)}$的直积$\Gamma^{(i)}\otimes\Gamma^{(j)}$的不可约表示含有$\text{A}_1$当且仅当$\Gamma^{(i)}$与$\Gamma^{(j)}$相同.
\end{theorem}
\subsection{群论在分子轨道理论中的应用}
\subsubsection{重叠积分}
根据前面的推导可知
\begin{theorem}[对称性相同的轨道才能有非零重叠]
    只有两个轨道生成相同的对称种类时,重叠积分$S$才可能不为$0$,并进而形成成键和反键轨道.
\end{theorem}
\subsubsection{对称性匹配线性组合}
构建对称性匹配线性组合(SALC)需要用到投影算符.
\begin{definition}[投影算符]
    对于群$G$中的不可约表示$\Gamma$,其投影算符$p^{\Gamma}$定义为
    \[p^{\Gamma}=\dfrac{1}{h}\sum_{g}\chi^{\Gamma}(g)g\]
\end{definition}
根据投影算符即可构建具有$\Gamma$对称性的轨道:
\[\psi^{\Gamma}=p^{\Gamma}\psi_i\]
\subsection{选律}
由波函数为$\psi_{\text{i}}$的始态向波函数为$\psi_{\text{f}}$的终态的跃迁产生的谱线强度决定于(电)跃迁偶极矩$\mu_{\text{fi}}$.这个矢量在$q$($q=x$, $y$, $z$)方向上的分量满足
\[\mu_{\text{q},\text{fi}}=-e\int_{\Omega}\psi_{\text{f}}^*r_q\psi_{\text{i}}\d\tau\]
跃迁偶极矩的分量是三个函数的乘积在全空间上的积分.类似地,考虑波函数$\psi_{\text{i}}$, $\psi_{\text{f}}$和$q$对应的不可约表示的直积是否包含全对称表示$\text{A}_1$即可.
\end{document}